\documentclass[11pt, a4paper]{article}

% Gemeinsame Style-Definitionen (TEXINPUTS muss templates/ enthalten — siehe scripts/build_tex.py)
% bfw-style.tex — gemeinsame Style-Definitionen für alle Handouts/Aufgabenhefte
% Voraussetzung: Kompilation mit xelatex (wegen fontspec)

\usepackage[ngerman]{babel}
% Babel-ngerman macht " zu einem aktiven Shorthand (z.B. für "a -> ä). In unseren
% Texten verwenden wir " als normales Zeichen — daher Shorthand komplett ausschalten.
\shorthandoff{"}
\usepackage{geometry}
\geometry{a4paper, top=2.4cm, bottom=2.4cm, left=2.3cm, right=2.3cm,
          headheight=15pt, headsep=0.7cm, footskip=1.1cm}

\usepackage{fontspec}
% Fallback-Kette: Source Sans Pro -> Calibri -> Segoe UI -> Latin Modern Sans
% (Latin Modern ist immer mit MiKTeX vorhanden)
\IfFontExistsTF{Source Sans Pro}{%
  \setmainfont{Source Sans Pro}[Ligatures=TeX]%
}{\IfFontExistsTF{Calibri}{%
  \setmainfont{Calibri}[Ligatures=TeX]%
}{\IfFontExistsTF{Segoe UI}{%
  \setmainfont{Segoe UI}[Ligatures=TeX]%
}{%
  \setmainfont{Latin Modern Sans}[Ligatures=TeX]%
}}}
\IfFontExistsTF{Source Code Pro}{%
  \setmonofont{Source Code Pro}[Scale=0.92]%
}{\IfFontExistsTF{Consolas}{%
  \setmonofont{Consolas}[Scale=0.92]%
}{%
  \setmonofont{Latin Modern Mono}[Scale=0.92]%
}}

% Gesperrte Versalien für Labels/Titelbalken (Kapitälchen-Ersatz, funktioniert
% mit jeder OpenType-Schrift — Latin Modern Sans hat keine echten Kapitälchen)
\newcommand{\bfwlabelfont}{\footnotesize\bfseries\addfontfeatures{LetterSpace=8.0}}

\usepackage{xcolor}
% Farbpalette: hell, freundlich, modern
\definecolor{bfwPrimary}{HTML}{0EA5A4}    % Petrol/Türkis - Primär
\definecolor{bfwPrimaryDark}{HTML}{0F766E} % dunkler Petrol für Headings
\definecolor{bfwAccent}{HTML}{F59E0B}     % Amber - Akzent
\definecolor{bfwSuccess}{HTML}{10B981}    % Grün - Erfolg / Lösung
\definecolor{bfwWarn}{HTML}{F97316}       % Orange - Achtung
\definecolor{bfwInk}{HTML}{0F172A}        % fast Schwarz
\definecolor{bfwInkSoft}{HTML}{475569}    % gedämpftes Grau
\definecolor{bfwBgSoft}{HTML}{F8FAFC}     % off-white
\definecolor{bfwBgTeal}{HTML}{ECFEFF}     % helles Türkis
\definecolor{bfwBgAmber}{HTML}{FEF3C7}    % helles Gelb
\definecolor{bfwBgRose}{HTML}{FFE4E6}     % helles Rosé für Eselsbrücken
\definecolor{bfwTabHead}{HTML}{E4F3F2}    % zarte Kopfzeilen-Hinterlegung für Tabellen

\usepackage[colorlinks=true, linkcolor=bfwPrimaryDark, urlcolor=bfwPrimaryDark]{hyperref}
\usepackage{lastpage}
\usepackage{enumitem}
\setlist[itemize]{leftmargin=*, itemsep=2pt, topsep=2pt}
\setlist[enumerate]{leftmargin=*, itemsep=2pt, topsep=2pt}

\usepackage[most]{tcolorbox}
\usepackage{booktabs}
\usepackage{colortbl}
\usepackage{tikz}
\usetikzlibrary{decorations.pathmorphing, positioning, shapes.geometric, arrows.meta}

% ---------------------------------------------------------------------------
% Einheitliches Tabellen-Design
%   - etwas mehr Zeilenluft, dezente graue Linien (wirkt auch mit \hline ruhiger)
%   - Kopfzeile einer Tabelle bei Bedarf zart hinterlegen: \rowcolor{bfwTabHead}
% ---------------------------------------------------------------------------
\renewcommand{\arraystretch}{1.25}
\arrayrulecolor{bfwInkSoft!50}
\setlength{\heavyrulewidth}{1pt}
\setlength{\lightrulewidth}{0.5pt}

% ---------------------------------------------------------------------------
% Boxen — klare farbige Titelbalken mit gesperrten Versal-Labels
% (Titel bleibt per Option überschreibbar: \begin{merksatz}[title={...}])
% ---------------------------------------------------------------------------
\tcbset{bfwbox/.style={%
  enhanced, breakable,
  boxrule=0.5pt, arc=2.5pt,
  left=10pt, right=10pt, top=8pt, bottom=8pt,
  toptitle=4pt, bottomtitle=4pt,
  titlerule=0pt,
  fonttitle=\bfwlabelfont,
  coltitle=white
}}

% Merkkästchen — wichtige Definition / Kernpunkt
\newtcolorbox{merksatz}[1][]{%
  bfwbox,
  colback=bfwBgTeal,
  colframe=bfwPrimaryDark,
  colbacktitle=bfwPrimaryDark,
  title={MERKE},
  #1
}

% Eselsbrücke — Mnemoniks
\newtcolorbox{eselsbruecke}[1][]{%
  bfwbox,
  colback=bfwBgRose,
  colframe=bfwAccent!85!black,
  colbacktitle=bfwAccent!85!black,
  title={ESELSBRÜCKE},
  #1
}

% Beispiel-Box
\newtcolorbox{beispiel}[1][]{%
  bfwbox,
  colback=bfwBgSoft,
  colframe=bfwInkSoft,
  colbacktitle=bfwInkSoft,
  title={BEISPIEL},
  #1
}

% Lösungs-Box (für Aufgabenhefte)
\newtcolorbox{loesung}[1][]{%
  bfwbox,
  colback=bfwSuccess!7,
  colframe=bfwSuccess!65!black,
  colbacktitle=bfwSuccess!65!black,
  title={LÖSUNG},
  #1
}

% Achtung-Box — typische Fehler / Stolperfallen
\newtcolorbox{achtung}[1][]{%
  bfwbox,
  colback=bfwWarn!6,
  colframe=bfwWarn!85!black,
  colbacktitle=bfwWarn!85!black,
  title={ACHTUNG},
  #1
}

% Formel-Box — für zentrale Formeln (groß, ruhig, ohne Titelbalken)
\newtcolorbox{formelbox}[1][]{%
  enhanced,
  colback=bfwBgTeal,
  colframe=bfwPrimaryDark,
  boxrule=1pt, arc=3pt,
  left=10pt, right=10pt, top=6pt, bottom=6pt,
  #1
}

% Glossar-Eintrag
\newcommand{\glossareintrag}[2]{%
  \par\noindent\textbf{\color{bfwPrimaryDark}#1} \enspace #2\par\smallskip
}

% ---------------------------------------------------------------------------
% Abschnitts-Design: farbiges Nummern-Quadrat + feine Linie unter \section,
% dezent farbige \subsection/\subsubsection
% ---------------------------------------------------------------------------
\usepackage{titlesec}
\newcommand{\bfwSecBadge}[1]{%
  \begin{tikzpicture}[baseline={([yshift=-0.62em]n.center)}]
    \node[fill=bfwPrimaryDark, text=white, font=\normalsize\bfseries,
          minimum width=1.55em, minimum height=1.55em, inner sep=1pt,
          rounded corners=1.5pt] (n) {#1};
  \end{tikzpicture}}
\titleformat{\section}
  {\Large\bfseries\color{bfwPrimaryDark}}
  {\bfwSecBadge{\thesection}}{0.6em}{}
  [\vspace{-0.2em}{\color{bfwPrimary!60}\titlerule[0.8pt]}]
\titleformat{\subsection}
  {\large\bfseries\color{bfwPrimaryDark}}
  {\textcolor{bfwPrimary}{\thesubsection}}{0.5em}{}
\titleformat{\subsubsection}
  {\normalsize\bfseries\color{bfwInk}}
  {\textcolor{bfwPrimary}{\thesubsubsection}}{0.4em}{}
\titlespacing*{\section}{0pt}{1.6em}{1em}
\titlespacing*{\subsection}{0pt}{1.2em}{0.5em}

% TikZ-Stil "skizzenhaft" — etwas weicher als Rough.js, aber stabil
\tikzset{
  roughbox/.style = {
    draw=bfwPrimaryDark, thick, fill=bfwBgTeal,
    rounded corners=4pt, inner sep=6pt
  },
  roughaccent/.style = {
    draw=bfwAccent, thick, fill=bfwBgAmber,
    rounded corners=4pt, inner sep=6pt
  },
  roughline/.style = {
    draw=bfwInkSoft, thick, -{Stealth[length=2.5mm]}
  }
}

% Bullet-Symbole (bewusst kein FontAwesome — vermeidet Font-Installations-Probleme)
\newcommand{\faLightbulb}{$\bullet$}
\newcommand{\faBookmark}{$\bullet$}

% ---------------------------------------------------------------------------
% Kopf-/Fußzeile
%   Kopf links:  Wortlogo „Wirtschaftsfachwirt"
%   Kopf rechts: Dokument-Kurztext, gesetzt via \kopfzeile{...}
%   Fuß links:   © Can Siebert · 2026 — Fuß rechts: Seite X von Y
%   Titelseite (Seite 1) bleibt ohne Kopf-/Fußzeile (\handouttitel setzt
%   \thispagestyle{empty}).
% ---------------------------------------------------------------------------
\usepackage{fancyhdr}
\newcommand{\bfwKopfzeileText}{}
\newcommand{\kopfzeile}[1]{\renewcommand{\bfwKopfzeileText}{#1}}
\pagestyle{fancy}
\fancyhf{}
\fancyhead[L]{\small\bfseries\color{bfwPrimaryDark}Wirtschaftsfachwirt}
\fancyhead[R]{\small\color{bfwInkSoft}\bfwKopfzeileText}
\renewcommand{\headrulewidth}{0.6pt}
\renewcommand{\headrule}{{\color{bfwPrimary}\hrule width\headwidth height 0.6pt}}
\fancyfoot[L]{\small\color{bfwInkSoft}\copyright\ Can Siebert\,\textperiodcentered\,2026}
\fancyfoot[R]{\small\color{bfwInkSoft}Seite \thepage\ von \pageref*{LastPage}}
% Auch \thispagestyle{plain} (z.B. durch \maketitle) soll leer bleiben:
\fancypagestyle{plain}{\fancyhf{}\renewcommand{\headrulewidth}{0pt}\renewcommand{\headrule}{}}

% ---------------------------------------------------------------------------
% \handouttitel{Obertitel}{Untertitel}{Kurs-Label}{Termin-Zeile}
%   Einheitlicher professioneller Titelblock: Dachzeile, farbige Headline,
%   Untertitel, farbige Trennlinie und dezente Meta-Zeile (Kurs/Termin/Dozent).
%   Ersetzt \title/\maketitle. Direkt nach \begin{document} aufrufen.
% ---------------------------------------------------------------------------
\newcommand{\handouttitel}[4]{%
  \thispagestyle{empty}%
  \begingroup
  \setlength{\parindent}{0pt}%
  {\bfwlabelfont\color{bfwPrimary}WIRTSCHAFTSFACHWIRT (IHK)\par}
  \vspace{0.55em}
  {\huge\bfseries\color{bfwPrimaryDark}#1\par}
  \vspace{0.5em}
  {\large\color{bfwInkSoft}#2\par}
  \vspace{0.9em}
  {\color{bfwPrimary}\rule{\linewidth}{2pt}}\par
  \vspace{0.55em}
  {\small\color{bfwInkSoft}#3\ \,\textperiodcentered\,\ #4\ \,\textperiodcentered\,\ Dozent: Can Siebert\par}
  \endgroup
  \vspace{1.4em}%
}


\kopfzeile{Handout BWL Lektion 1 \textperiodcentered{} Teil 2}

\begin{document}
\handouttitel{Lektion~1 \textperiodcentered{} Teil~2: Absatz \& Marketing}%
  {Vom Marketingbegriff über Marktforschung und Strategie zum Marketing-Mix --- inkl. Produktlebenszyklus und Break-even}%
  {1.2 Betriebliche Funktionen}%
  {Sa, 24.10.2026 \textperiodcentered{} Session 2}

\begin{merksatz}[title={\bfseries Worum geht es in Teil 2?}]
In Teil~1 haben wir das Zielsystem, die Produktion und die Logistik behandelt --- nun
schließen wir den Wertschöpfungsprozess mit der Funktion \textbf{Absatz/Marketing} ab:
der marktorientierten Verwertung der erstellten Leistungen. Wir klären den modernen
\textbf{Marketingbegriff} (Kundenorientierung), die Methoden der \textbf{Marktforschung},
\textbf{Marketingziele und -strategien} (Ansoff-Matrix) und die vier Instrumente des
\textbf{Marketing-Mix (4P)} --- inklusive \textbf{Produktlebenszyklus} und
\textbf{Break-even-Rechnung} am Fallbeispiel. Die übrigen Funktionen (Rechnungswesen,
Finanzierung, Controlling, Personal) folgen in Lektion~2.
\end{merksatz}

\tableofcontents
\vspace{1em}
\hrule
\vspace{1em}

\section{Absatz und Marketing}

\subsection{Der Marketingbegriff}

Das \textbf{Marketingdenken} von heute wird dominiert durch die Konzentration auf den
\textbf{Kunden}: Die Leistung eines Unternehmens besteht in der \emph{Problemlösung für
den potenziellen Käufer}. Marketing ist damit mehr als Werbung oder Verkauf --- es ist
die konsequent \textbf{markt- und kundenorientierte Führung} des gesamten Unternehmens.
Alle betrieblichen Funktionen sind auf den Markt auszurichten. Einsatzfelder des
Marketings sind das Konsumgüter-, Investitionsgüter- und Dienstleistungsmarketing, das
internationale Marketing sowie das Marketing von \emph{Non-Profit-Organisationen} ---
auch dort ist Marketing notwendig. Das Investitionsgütermarketing setzt am
\emph{abgeleiteten} Bedarf an (die Nachfrage nach Maschinen folgt aus der Nachfrage nach
Konsumgütern).

Das \textbf{Marketingmanagement} verläuft als Kreislauf: \emph{Information/Analyse}
(Markt, Zielgruppe, Absatzwege) $\rightarrow$ \emph{Zielsetzung} $\rightarrow$
\emph{alternative Strategien} $\rightarrow$ \emph{optimale Strategie} $\rightarrow$
\emph{Voraussetzungen schaffen} $\rightarrow$ \emph{Realisierung} (Marketing-Mix)
$\rightarrow$ \emph{Controlling} --- mit Rückkopplungen auf allen Stufen.

\begin{merksatz}[title={\bfseries Prüfungsfalle}]
Ausgangspunkt aller Marketingaktivitäten sind \textbf{Informationen über den Markt}
(Zielgruppe, Absatzwege) --- \emph{nicht} das fertige Produkt. Diese
Richtig-falsch-Frage taucht im Übungsband und in Prüfungen regelmäßig auf.
\end{merksatz}

\subsection{Marktforschung}

Die \textbf{Marktforschung} beschafft und analysiert systematisch Informationen über
Märkte (Kunden, Wettbewerber, Absatzmittler) als Grundlage der Marketingentscheidungen:

\begin{itemize}
  \item \textbf{Primärforschung} (field research): Erhebung \emph{neuer} Daten ---
    \textbf{Befragung} (schriftlich, telefonisch, online, Interview),
    \textbf{Beobachtung} (z.\,B. Kundenlaufstudien), \textbf{Experiment} (Markt-/Storetest)
    und \textbf{Panel} (wiederholte Erhebung bei gleichbleibendem Teilnehmerkreis).
  \item \textbf{Sekundärforschung} (desk research): Auswertung \emph{vorhandener} Daten
    (interne Statistiken, amtliche Statistik, Verbandszahlen, Studien) --- schneller und
    günstiger, aber oft nicht passgenau und veraltet.
\end{itemize}

Zu unterscheiden sind \textbf{Marktanalyse} (Momentaufnahme des Marktes) und
\textbf{Marktbeobachtung} (Entwicklung im Zeitverlauf). Aus Kostengründen gilt: erst
Sekundär-, dann ergänzend Primärforschung.

\subsection{Marketingziele}

Marketingziele lassen sich in zwei Gruppen unterscheiden:

\begin{itemize}
  \item \textbf{Ökonomische Ziele} --- messbare, wirtschaftliche Größen: Umsatz, Absatz,
    Gewinn, Marktanteil, Deckungsbeitrag, Distributionsgrad.
  \item \textbf{Psychografische Ziele} --- nicht direkt messbare, im Kopf des Kunden
    verankerte Größen: Bekanntheitsgrad, Image, Kundenzufriedenheit, Kundenbindung,
    Kaufpräferenzen und Wiederkaufabsicht.
\end{itemize}

\begin{eselsbruecke}
\textbf{Öko = Zahlen, Psycho = Köpfe}: Ökonomische Ziele stehen in der
Ergebnisrechnung, psychografische Ziele im Kopf des Kunden --- sie wirken aber
langfristig auf die ökonomischen Ziele.
\end{eselsbruecke}

\subsection{Marketingstrategien}

Auf Basis von Marktsegmentierung (Aufteilung des Gesamtmarktes in homogene
Käufergruppen nach demografischen, sozioökonomischen und psychografischen Kriterien)
werden Strategien festgelegt, z.\,B. nach der \textbf{Produkt-Markt-Matrix (Ansoff)}:

\begin{center}
\begin{tabular}{p{3.4cm} p{5.2cm} p{5.2cm}}
 & \textbf{gegenwärtige Produkte} & \textbf{neue Produkte}\\
\hline
\textbf{gegenwärtige Märkte} & Marktdurchdringung (Marktanteil ausbauen) & Produktentwicklung (Innovationen für Stammkunden)\\
\textbf{neue Märkte} & Marktentwicklung (neue Regionen, neue Zielgruppen) & Diversifikation (neue Produkte auf neuen Märkten)\\
\end{tabular}
\end{center}

Weitere Grundentscheidungen: \textbf{Massenmarktstrategie} vs.
\textbf{Segmentierungsstrategie} (Nischen), \textbf{Präferenzstrategie} (Marke, Qualität,
höherer Preis) vs. \textbf{Preis-Mengen-Strategie} (niedriger Preis, große Menge).

\subsection{Der Marketing-Mix (4P)}

Der \textbf{Marketing-Mix} ist das Ergebnis der strategischen Marketingplanung: der
aufeinander abgestimmte Einsatz der vier Instrumentalbereiche (englisch: Product, Price,
Place, Promotion).

\begin{center}
\begin{tikzpicture}[font=\footnotesize,
  p/.style={roughbox, minimum width=5.9cm, minimum height=1.85cm, align=center, inner sep=5pt},
  mitte/.style={roughaccent, ellipse, minimum width=3.6cm, minimum height=1.5cm, align=center, font=\small},
  lin/.style={draw=bfwInkSoft, thick}]
  \node[mitte] (m) at (7,-1.7) {\textbf{Marketing-Mix}\\{\scriptsize Zielgruppe/Kunde}};
  \node[p] (prod)  at (2.9,0)    {\textbf{Product --- Produktpolitik}\\{\scriptsize Qualität, Design, Sortiment,}\\{\scriptsize Marke, Service/Kundendienst}};
  \node[p] (price) at (11.1,0)   {\textbf{Price --- Preis-/Kontrahierungspolitik}\\{\scriptsize Preis, Rabatte, Boni, Skonti,}\\{\scriptsize Liefer- und Zahlungsbedingungen}};
  \node[p] (place) at (2.9,-3.4) {\textbf{Place --- Distributionspolitik}\\{\scriptsize direkte/indirekte Absatzwege,}\\{\scriptsize Marketinglogistik}};
  \node[p] (promo) at (11.1,-3.4){\textbf{Promotion --- Kommunikationspolitik}\\{\scriptsize Werbung, Verkaufsförderung, PR,}\\{\scriptsize persönlicher Verkauf, Social Media}};
  \draw[lin] (m) -- (prod); \draw[lin] (m) -- (price);
  \draw[lin] (m) -- (place); \draw[lin] (m) -- (promo);
\end{tikzpicture}\\[0.3em]
{\small\itshape Die 4P des Marketing-Mix --- alle vier Instrumente müssen aufeinander
und auf die Zielgruppe abgestimmt sein.}
\end{center}

\begin{itemize}
  \item \textbf{Produkt- und Leistungspolitik (Product):} Produktgestaltung (Qualität,
    Design, Verpackung), Sortiments-/Programmpolitik, Produktinnovation, -variation,
    -elimination, Markenpolitik sowie \textbf{Service/Kundendienst} --- vor dem Kauf
    (Information, technische Beratung, Problemlösungsvorschläge, Finanzierungsangebote),
    in der Kaufphase (Lieferung, Montage, Umtauschrecht) und nach dem Kauf
    (Garantieleistungen, Wartung und Reparatur, Ersatzteilversorgung).
  \item \textbf{Preis- und Kontrahierungspolitik (Price):} Preisfestsetzung
    (kosten-, nachfrage-, konkurrenzorientiert), Preisdifferenzierung, Rabatte, Boni und
    Skonti, Liefer- und Zahlungsbedingungen, Absatzkredite.
  \item \textbf{Distributionspolitik (Place):} Wahl der Absatzwege --- direkter Absatz
    (eigener Vertrieb, Onlineshop, Werksverkauf) oder indirekter Absatz (Groß- und
    Einzelhandel, Handelsvertreter) --- sowie die physische Distribution
    (Marketinglogistik): das Produkt zum richtigen Zeitpunkt in der gewünschten Menge an
    den Ort der Nachfrage bringen.
  \item \textbf{Kommunikationspolitik (Promotion):} \textbf{Werbung} (bezahlte,
    produktbezogene Kommunikation zur Absatzförderung), \textbf{Verkaufsförderung}
    (Sales Promotion: Aktionen am Point of Sale), \textbf{Public Relations}
    (Öffentlichkeitsarbeit: Vertrauen und gutes Image für das \emph{Unternehmen} ---
    nicht unmittelbar produktbezogen), persönlicher Verkauf, Sponsoring, Messen,
    Online-/Social-Media-Marketing.
\end{itemize}

\begin{eselsbruecke}
Stufen der Werbewirkung --- \textbf{AIDA}: \textbf{A}ttention (Aufmerksamkeit wecken),
\textbf{I}nterest (Interesse erzeugen), \textbf{D}esire (Kaufwunsch auslösen),
\textbf{A}ction (Kaufhandlung herbeiführen). Ein Klassiker für Fragen zur
Kommunikationspolitik.
\end{eselsbruecke}

\subsection{Der Produktlebenszyklus}

Grundlage vieler produktpolitischer Entscheidungen ist der \textbf{Produktlebenszyklus}
mit den Phasen \emph{Einführung $\rightarrow$ Wachstum $\rightarrow$ Reife $\rightarrow$
Sättigung $\rightarrow$ Degeneration (Rückgang)}: In der Einführungsphase entstehen
meist Verluste (hohe Entwicklungs- und Markteinführungskosten bei erst geringen
Umsätzen), das Umsatzmaximum liegt in der Sättigungsphase --- danach sind
\textbf{Relaunch} (Neubelebung), \textbf{Variation} oder \textbf{Elimination} angezeigt.

\begin{center}
\begin{tikzpicture}[x=1.8cm, y=0.85cm, font=\scriptsize]
  % Achsen
  \draw[-{Stealth}, thick, bfwInkSoft] (0,-1.4) -- (0,3.9) node[above, font=\footnotesize] {Umsatz / Gewinn};
  \draw[-{Stealth}, thick, bfwInkSoft] (0,0) -- (7.4,0) node[right, font=\footnotesize] {Zeit};
  % Phasengrenzen
  \foreach \x in {1.4,2.8,4.2,5.6} \draw[dashed, bfwInkSoft!50] (\x,-1.2) -- (\x,3.4);
  \node at (0.7,3.62)  {Einführung};
  \node at (2.1,3.62)  {Wachstum};
  \node at (3.5,3.62)  {Reife};
  \node at (4.9,3.62)  {Sättigung};
  \node at (6.5,3.62)  {Degeneration};
  % Umsatzkurve
  \draw[very thick, bfwPrimaryDark] plot [smooth, tension=0.7] coordinates
    {(0,0) (1.4,0.7) (2.8,2.2) (4.2,2.95) (5.0,3.05) (5.6,2.9) (7,1.5)};
  \node[bfwPrimaryDark, font=\footnotesize\bfseries] at (6.65,2.45) {Umsatz};
  % Gewinnkurve
  \draw[very thick, dashed, bfwAccent!85!black] plot [smooth, tension=0.7] coordinates
    {(0,-1.0) (1.4,-0.35) (2.0,0) (2.8,0.8) (4.2,1.35) (5.6,0.9) (7,0.1)};
  \node[bfwAccent!85!black, font=\footnotesize\bfseries] at (5.15,1.85) {Gewinn};
  \node[bfwInkSoft, align=left] at (2.6,-1.0) {Verlustzone: hohe Entwicklungs- und Einführungskosten};
\end{tikzpicture}\\[0.3em]
{\small\itshape Der Produktlebenszyklus: Der Gewinn wird erst in der Wachstumsphase
positiv und sinkt vor dem Umsatz --- am Ende stehen Relaunch, Variation oder
Elimination.}
\end{center}

\begin{beispiel}
Die Nordwest Küchengeräte AG bringt einen neuen Mixer auf den Markt (Fixkosten
150.000~€, variable Kosten 14~€/Stück, Marktpreis 30~€, erwartete Absatzmenge
150.000 Stück). Erfolgskontrolle: Gewinn $= 150.000 \cdot (30 - 14) - 150.000 =
2.400.000 - 150.000 = 2.250.000$~€. Die Gewinnschwelle (Break-even) liegt bei
$150.000 / 16 = 9.375$ Stück --- das Projekt ist unter Kosten- und
Erfolgsgesichtspunkten zu realisieren (Rechenweg im Detail: Abschnitt~2).
\end{beispiel}

\section{Formeln \& Rechenbeispiele}

Die prüfungsrelevante Berechnung dieser Session mit Formel, Zahlenbeispiel und
Schritt-für-Schritt-Rechenweg. (Produktivität, Wirtschaftlichkeit und Lagerkennzahlen:
siehe Handout zu Teil~1.)

\subsection{Break-even-Menge (Gewinnschwelle)}

\begin{merksatz}[title={\bfseries Formel: Break-even-Menge}]
\[
\text{Break-even-Menge} \;=\; \frac{\text{Fixkosten}}{\text{Verkaufspreis} - \text{variable Stückkosten}}
\]
Der Nenner ist der \textbf{Stückdeckungsbeitrag}. An der Gewinnschwelle gilt:
Erlöse $=$ Gesamtkosten, der Gewinn ist \textbf{null}. Ergänzend:
Gewinn $=$ Menge $\times$ (Preis $-$ variable Stückkosten) $-$ Fixkosten.
\end{merksatz}

\begin{beispiel}
Mixer der Nordwest Küchengeräte AG: Fixkosten 150.000~€, Verkaufspreis 30~€,
variable Kosten 14~€ je Stück.
\begin{enumerate}
  \item Stückdeckungsbeitrag berechnen: $30 - 14 = 16$~€ je Stück.
  \item Fixkosten durch Stückdeckungsbeitrag teilen:
    $150.000 \div 16 = 9.375$ \textbf{Stück}.
  \item Interpretieren: Ab dem 9.376. Stück macht das Produkt Gewinn --- bei der
    erwarteten Absatzmenge von 150.000 Stück beträgt er
    $150.000 \times 16 - 150.000 = 2.250.000$~€.
\end{enumerate}
\end{beispiel}

\begin{center}
\begin{tikzpicture}[x=1.05cm, y=0.72cm, font=\scriptsize]
  % Achsen
  \draw[-{Stealth}, thick, bfwInkSoft] (0,0) -- (0,6.7) node[above, font=\footnotesize] {€};
  \draw[-{Stealth}, thick, bfwInkSoft] (0,0) -- (12.6,0) node[right, font=\footnotesize] {Menge (Stück)};
  % Fixkosten
  \draw[thick, bfwInkSoft] (0,1.8) -- (12,1.8) node[pos=0.94, below] {Fixkosten};
  % Gesamtkosten
  \draw[very thick, bfwAccent!85!black] (0,1.8) -- (12,5.0)
    node[pos=0.38, above, sloped] {Gesamtkosten (Fixkosten $+$ variable Kosten)};
  % Erlöse
  \draw[very thick, bfwPrimaryDark] (0,0) -- (12,6.3)
    node[pos=0.97, above, sloped] {Erlöse};
  % Break-even-Punkt
  \draw[dashed, bfwInkSoft] (6.97,0) -- (6.97,3.66);
  \fill[bfwAccent] (6.97,3.66) circle (3pt);
  \node[align=center, fill=bfwBgAmber, draw=bfwAccent!85!black, rounded corners=3pt,
        inner sep=3pt] (bep) at (3.6,5.7)
    {\textbf{Break-even-Punkt (Gewinnschwelle)}\\Erlöse $=$ Gesamtkosten, Gewinn $=$ 0\\hier: 9.375 Stück (Mixer-Beispiel)};
  \draw[roughline] (bep.south east) -- (6.8,3.9);
  \node[bfwInkSoft, align=center] at (3.4,2.25) {Verlustzone};
  \node[bfwInkSoft, align=center] at (10.2,4.95) {Gewinnzone};
\end{tikzpicture}\\[0.3em]
{\small\itshape Break-even-Diagramm: Links vom Schnittpunkt liegen die Gesamtkosten
über den Erlösen (Verlust), rechts davon entsteht Gewinn.}
\end{center}

\subsection{Deckungsbeitrag und Zielgewinn-Menge}

\begin{merksatz}[title={\bfseries Formeln: Deckungsbeitrag und Zielgewinn-Menge}]
\[
\text{Stückdeckungsbeitrag } db \;=\; \text{Verkaufspreis} - \text{variable Stückkosten}
\]
\[
\text{Gesamtdeckungsbeitrag} \;=\; db \times \text{Menge}
\quad\;
\text{Menge für einen Zielgewinn} \;=\; \frac{\text{Fixkosten} + \text{Zielgewinn}}{db}
\]
Der Deckungsbeitrag „deckt`` zunächst die Fixkosten --- erst wenn diese vollständig
gedeckt sind (Break-even), entsteht Gewinn. Die Zielgewinn-Formel ist nur die
Break-even-Formel mit einem Aufschlag im Zähler.
\end{merksatz}

\begin{beispiel}
Mixer der Nordwest Küchengeräte AG (Preis 30~€, variable Kosten 14~€, Fixkosten
150.000~€). Die Geschäftsleitung fordert einen \emph{Zielgewinn von 250.000~€} im
ersten Jahr.
\begin{enumerate}
  \item Stückdeckungsbeitrag: $30 - 14 = 16$~€ je Stück.
  \item Benötigte Menge: $(150.000 + 250.000) \div 16 = 25.000$ \textbf{Stück}.
  \item Probe über den Gesamtdeckungsbeitrag: $25.000 \times 16 = 400.000$~€ ---
    davon 150.000~€ für die Fixkosten, es verbleiben genau 250.000~€ Gewinn.
  \item Interpretieren: Bei der erwarteten Absatzmenge von 150.000 Stück ist der
    Zielgewinn realistisch --- schon 25.000 verkaufte Mixer genügen.
\end{enumerate}
\end{beispiel}

\begin{merksatz}[title={\bfseries Wichtig für die Prüfungsvorbereitung}]
Sie sollten aus Teil~2 sicher können: den \emph{Marketingbegriff} (Kundenorientierung,
Ausgangspunkt Marktinformationen), die Unterscheidung \emph{Primär-/Sekundärforschung}
mit ihren Methoden, \emph{ökonomische vs. psychografische Marketingziele}, die vier
Strategien der \emph{Ansoff-Matrix}, den \emph{Marketing-Mix} mit den
Entscheidungsebenen der vier Instrumente, die Phasen des \emph{Produktlebenszyklus}
sowie die Formeln der \emph{Break-even-Menge}, des \emph{Deckungsbeitrags} und der
\emph{Zielgewinn-Menge} (Abschnitt~2). Üben Sie mit dem
Aufgabenheft --- z.\,B. der Marketing-Mix-Fallaufgabe --- und dem Quiz zu Teil~2!
\end{merksatz}

\vspace{1em}
\hrule
\vspace{0.8em}
\noindent\textsf{\small Quellen: Rahmenplan Wirtschaftsbezogene Qualifikationen (DIHK), Abschnitt 1.2.1;
Übungsband Betriebswirtschaft (DIHK-Bildungs-GmbH), Kap.~1. Der Auszug aus dem Textband
Betriebswirtschaft wird nachgereicht. Nächste Lektion (Sa, 07.11.2026): Betriebliche Funktionen II ---
Rechnungswesen, Finanzierung \& Investition, Controlling, Personalwirtschaft.}

\end{document}
